
\section{Introduction}\label{sec:intro}

Methods for learning causal structure from observational data are mostly validated on
simulated data, where the generating mechanism is known and the assumptions of the method
can be matched to it. Real measurements differ in ways that simulation rarely reproduces:
mechanisms are nonlinear and heterogeneous, sensors are collinear, unobserved processes
induce dependence between measured variables, and rows are often not independent draws.
Benchmarks with real data and a known structure exist, but each covers a narrow slice:
the Tübingen cause--effect pairs \citep{mooij2016distinguishing}, the protein-signalling
data of \citet{sachs2005causal}, whose reference network is a contested consensus
\citep{mooij2020joint}, and, more recently, the Causal Chambers, physical devices whose
graph is known from their construction \citep{gamella2025chamber}.

This paper asks a plain question: when methods are applied unchanged to several real
data families, which ones work, and do recent ideas beat simple baselines? We evaluate
fourteen methods on four families of real data, under a protocol in which every method and
every setting was fixed and recorded before the data it is scored on were run, with the
exceptions listed in Section~\ref{sec:provenance}. The methods are eight established
baselines (PC, GES, DirectLiNGAM, NOTEARS for graphs; IGCI, RECI, ANM-HSIC and a pairwise
LiNGAM rule for pairs) and six methods that we developed during the project that produced
this study: sequential certificates for adjacency and direction, PC with a nonparametric
test, a block-replicated additive method, and two classifiers trained on simulated data.
Including our own methods is deliberate. Several were designed to address failures of the
baselines, and the comparison tells whether the design worked on measurements.

\paragraph{Contributions.}
\begin{enumerate}
\item A multi-dataset evaluation with a recorded protocol. It covers $102$ Tübingen
  pairs, $216$ pair analyses of edges of the Causal Chambers, the Sachs data, and
  $21$ chamber recordings with a known graph, and it reports a strict score against the
  reference graph and a score adjusted for dependences the graph cannot explain.
\item The result that, on pairs, a regression-error rule is the best or tied-best method on
  every real data set, and a classifier trained on simulated pairs, accurate on its
  simulated test set, is at chance on Tübingen (Section~\ref{sec:pairs}).
\item The result that, on graphs, the ranking of methods depends on the data family: DirectLiNGAM
  and GES on chamber recordings with independent rows, methods with nonparametric tests on the
  Sachs data, and no method on the random-walk recordings (Section~\ref{sec:graphs}).
\item A diagnosis of why the walk recordings defeat every method: their exogenous variables are
  random walks with an effective sample size of a few observations, so tests that assume
  independent rows reject graph-implied independences about five times more often than a block
  test, and a reference set built with them is contaminated (Section~\ref{sec:serial}).
\item The places where our own methods do help: control of false edges under repeated
  inspection and under nonlinear mechanisms in simulation (Section~\ref{sec:monitoring}).
\end{enumerate}

\paragraph{What the paper does not claim.} It does not claim that the methods we did not
evaluate (LOCI, QCCD, CAM, DAGMA, PCMCI+ and amortised methods such as AVICI
\citep{lorch2022amortized}) fail on these data, and it does not claim that the data
sets are representative of causal discovery in general. Several data sets have few
reference edges, so most differences between methods are within sampling error, and we
give intervals where we can.
